%!TEX root = ../thesis.tex

本研究では，自然言語の指示のみで画像から対象を切り出すゼロショットセグメンテーションに物体追跡を追加し，リアルタイム性を向上する手法を提案した．提案手法の評価では，処理速度を2.7倍に向上できた一方で，IoUが0.107低下するトレードオフも見られた．

また，実環境への適用可能性を調査するため，提案手法をAI-Formulaに応用した．その結果，モビリティは指定したコースを逸脱せず，障害物の手前で停止した．さらに，システム構成を改良して，大学のキャンパス内での走行実験を行った．その結果，モビリティは被験者を追従し，従来手法よりも処理速度が向上した．

今後は，実験の試行回数を増加したり，データの解析をしたり，定量的な評価に取り組んでいく．また，さらなる高速化に取り組み，AI-Formulaのレースや他のロボット技術チャレンジへの応用を目指す．

\newpage
